\section{Topology}

\begin{exercise}[A phase outside the solenoid]
\textbf{Experiment.}
Electron interference fringes shift when magnetic flux through an excluded
solenoid changes, even though the magnetic field vanishes along both electron
paths \cite{aharonov-bohm,tonomura}.  Let \(X\) be the field-free region and
let \(A\) be a closed one-form there.

Show that the relative phase is
\[
        \exp\!\left(\frac{iq}{\hbar}\oint_\gamma A\right),
\]
where \(\gamma\) is the loop formed by the two paths.  Prove invariance under
\(A\mapsto A+df\).  Show that the phase depends only on the de Rham class of
\(A\) and the homology class of \(\gamma\).  Compute it for the complement of
an ideal solenoid, and explain why the observed fringe shift detects
holonomy rather than a local magnetic force on the electrons.
\end{exercise}

\begin{exercise}[A geometric neutron phase]
\textbf{Experiment.}
Polarized neutrons transported through a cyclic sequence of magnetic fields
acquire a phase proportional to the solid angle enclosed by the field
direction on the sphere \cite{berry-neutron}.  Let \(L\to S^2\) be the line
whose fiber is the instantaneous positive-spin eigenspace.

Use orthogonal projection of the trivial connection to construct a connection
on \(L\).  Calculate its curvature and show that the holonomy around a loop is
half the enclosed solid angle, modulo \(2\pi\).  Explain why changing the phase
of each chosen eigenvector changes the local connection form but not the
measured interference phase.  Determine why no globally smooth nonvanishing
choice of positive-spin eigenvector exists.
\end{exercise}

\begin{exercise}[Integer Hall plateaus]
\textbf{Experiment.}
At low temperature, a two-dimensional electron system displays Hall
conductance locked to integer multiples of \(e^2/h\), with plateaus stable
under disorder and small changes of magnetic field \cite{von-klitzing}.
For a gapped periodic model, let the occupied states form a vector bundle
\(E\) over the Brillouin torus.

Construct the projected connection on \(E\) and its curvature \(F\).  Show
that
\[
        \frac{i}{2\pi}\int_{\mathbb T^2}\Tr(F)\in\Z.
\]
Use the response formula to identify this integer with Hall conductance
\cite{tknn}.  Prove that it is unchanged under a continuous deformation that
keeps the spectral gap open.  Explain both the plateau's stability and the
need for a gap closing when the measured integer changes.
\end{exercise}

\begin{exercise}[One particle per pump cycle]
\textbf{Experiment.}
An adiabatically driven one-dimensional lattice can transport an integer
number of particles per cycle; cold-atom measurements observe this quantized
pumping \cite{lohse}.  Treat crystal momentum and drive time as coordinates on
a torus, and let \(E\) be the occupied-state bundle over that torus.

Relate the transported charge in one period to the curvature of the projected
connection on \(E\).  Show that the result is its first Chern number and is
therefore integral.  Explain why small errors in the drive do not change the
transported integer while adiabaticity and the energy gap persist.  Identify
what the experiment would show if the path crossed a degeneracy.
\end{exercise}

\begin{exercise}[Zero modes at a boundary]
\textbf{Experiment.}
Topological materials support robust boundary modes whose net direction of
propagation is tied to a bulk integer.  Let \(D\) be the Fredholm operator
obtained from the boundary problem, with finite-dimensional kernel and
cokernel.

Show that
\[
        \operatorname{ind}D=\dim\ker D-\dim\operatorname{coker}D
\]
is unchanged by sufficiently small perturbations.  Relate this analytic
integer to the bulk Chern number through an index pairing.  Explain why
disorder may deform individual boundary states but cannot change the net
number of chiral channels without closing the bulk gap or mixing with an
opposite boundary.  State the experimentally testable transport consequence
of this equality.
\end{exercise}

\begin{exercise}[Quantized superconducting flux]
\textbf{Experiment.}
Magnetic flux trapped in a superconducting ring occurs in integer multiples
of \(h/2e\).  Let the nonzero condensate define a phase on the ring and let
\(A\) be the electromagnetic connection.

Show that single-valuedness makes the phase winding an integer.  Combine its
differential with \(A\) to form the gauge-invariant circulation and recover
the observed flux unit when the bulk current vanishes.  Explain why local
changes of condensate phase do not alter the integer.  Identify the event
required for the ring to move between adjacent flux sectors.
\end{exercise}

\begin{exercise}[Weak localization in a ring]
\textbf{Experiment.}
The conductance of a small disordered metallic ring oscillates as magnetic
flux through the ring is varied, even when electrons travel diffusively.
Pair each closed diffusive path with its time reverse.

Show that the two amplitudes have equal dynamical phase but acquire opposite
magnetic holonomies.  Derive the flux period of their interference correction.
Explain why averaging over disorder suppresses many path pairs but preserves
the time-reversed pair.  Interpret the oscillation as sensitivity to the
connection's holonomy rather than to a force along the wire.
\end{exercise}

\begin{exercise}[Graphene's half-integer offset]
\textbf{Experiment.}
Graphene exhibits a quantum Hall sequence shifted relative to that of an
ordinary parabolic band.  Near a band-touching point, the low-energy
eigenspaces form lines over a loop in momentum space.

Transport an eigenline around the loop and compute its geometric phase.
Explain how the resulting sign changes the semiclassical quantization rule and
produces the observed Hall offset.  Show that the phase cannot be removed by
choosing eigenvectors smoothly along the entire loop while respecting
single-valuedness.  State what perturbation can gap the touching point and
alter the inference.
\end{exercise}

\begin{exercise}[A topological interface mode]
\textbf{Experiment.}
Engineered one-dimensional lattices with alternating couplings display a
localized mode where the coupling pattern switches order.  Let the gapped
bulk Hamiltonian possess a symmetry which pairs positive and negative
energies.

Associate a winding integer to each translation-invariant bulk.  Show that
the two coupling orders have different integers.  Explain why an interface
must then support a midgap state and why its wave function is localized.
Determine which perturbations can move or remove the state, and relate the
answer to whether the protecting symmetry and bulk gap survive.
\end{exercise}

\begin{exercise}[Josephson voltage oscillations]
\textbf{Experiment.}
A constant voltage across a Josephson junction produces an alternating
supercurrent whose frequency obeys \(f=2eV/h\).  Let the two superconductors
carry condensate phases connected by the electromagnetic gauge field.

Construct the gauge-invariant phase difference across the junction.  Use its
time evolution to derive the measured frequency--voltage relation.  Explain
why the coefficient reveals charge \(2e\), why a local phase convention
cannot alter the prediction, and how phase winding makes the oscillation a
global comparison rather than an observable absolute phase.
\end{exercise}
